
Sequential hypothesis testing requires two distinct decisions: \textit{when} to stop
collecting data and \textit{what} to conclude once stopped. Conflating these steps,
as in the HDI+ROPE algorithm, risks stopping on imprecise, unrepresentative posteriors
and produces a systematic false-positive rate ($\sim\!6\%$ in our simulations, even
when $\theta_{\rm true}=\theta_{\rm null}$).
\cite{kruschke2015doing}'s Precision is the Goal (PitG) eliminates this bias by
decoupling the two steps: it stops only once the posterior HDI narrows below a
pre-specified precision goal $\omega_{\rm goal}$, then applies the HDI+ROPE decision
rule. In the fair-coin case, PitG achieves zero false positives, but at the cost of a
strikingly high inconclusive rate of 63.3\%, because precision and conclusiveness
are not required to coincide.
Moreover, unlike NHST, which can only reject or fail to reject the null,
the HDI+ROPE decision rule enables positive \textit{acceptance} of the null
when the entire HDI falls within the ROPE, a critical advantage for confirmatory
research areas such as bioequivalence, safety testing, and software regression testing
(verifying that a new deployment has not degraded existing behaviour known as do-no-harm).

We have proposed Enhanced Precision is the Goal (ePitG), which re-couples the two
criteria conservatively: data collection continues until the HDI is simultaneously
narrower than $\omega_{\rm goal}$ \textit{and} falls cleanly inside or outside the
ROPE\@. This single change transforms the picture. At $\theta_{\rm true}=\theta_{\rm null}=0.5$,
ePitG achieves a 97.7\% conclusiveness rate with zero false positives, compared to
PitG's 36.7\% --- a $2.7\times$ gain in conclusiveness at a median cost of only 4.7\%
more samples. The efficiency gain is most pronounced near the ROPE boundaries, precisely
where correct decisions matter most; once $\theta_{\rm true}$ moves well outside the
ROPE, ePitG converges to PitG and both stop at $N_{\rm goal}(\theta_{\rm true})$.
These figures correspond to $\omega_{\rm goal}=0.08$, i.e.\ a $\pm 4$ percentage
point margin of error, a standard reporting precision in single-group binomial analyses
such as opinion polling and survey-based research. The qualitative advantage of ePitG over PitG is robust
across precision goals: looser goals (larger $\omega_{\rm goal}$) worsen PitG's
inconclusive rate and amplify ePitG's gain, while tighter goals improve PitG's
conclusiveness at the cost of proportionally larger samples for both methods.

Both precision-based methods share an important analytical property: the expected stop
iteration is given in closed form by
$N_{\rm goal}(\theta,\omega_{\rm goal}) \propto \theta(1-\theta)/\omega_{\rm goal}^2$
(Equation~\ref{eq:pitg_stop_iteration}), enabling prospective planning before any data
are collected. $N_{\rm goal}(\theta_{\rm null})$ is always a theoretical minimum
for the required sample size. Because $\theta(1-\theta)$ is maximised at $\theta=0.5$,
this minimum is largest for the fair-coin null, making it a safe worst-case planning
ceiling. For other nulls, $\theta_{\rm true}$ may lie closer to $0.5$ than
$\theta_{\rm null}$ does, raising the actual minimum above the planned one;
$N_{\rm max}$ should be set generously to cover this (see Discussion).

The framework extends naturally to single-group continuous data and two-group
comparisons (A/B testing); the adaptations and sample-size planning equations are
detailed in Appendix~\ref{app:extensions}, and a full empirical evaluation across
these settings is left for future work.
To support adoption, all algorithms are implemented in openly available
code\footnote{\textcolor{red}{\textbf{[TODO: insert repo URL]}}
\href{https://github.com/placeholder/precision-goal}{https://github.com/placeholder/precision-goal}},
and an interactive online calculator for prospective sample-size planning via
Equation~\ref{eq:pitg_stop_iteration} is available at
\href{https://r-w-t-y.streamlit.app}{https://r-w-t-y.streamlit.app}.
